\documentclass[fleqn]{article}
\usepackage{amsmath}
\usepackage{amsfonts}
\usepackage{enumitem}
\title{Solución Problemas Números Complejos}
\author{Juan Rodríguez}

\setlength{\parindent}{0pt}

\begin{document}
	\maketitle
	\section*{Apuntes}
	\begin{itemize}
		\item $i^{2} = -1$
		\item $i^{3} = i^{2} \cdot i = -i$
		\item $i^{4} = i^{2} \cdot i^{2} = 1$
		\item $(a+bi) \cdot (c+di) = (ac-bd+adi+bci)$
		\item $\frac{a+bi}{c+di} = \frac{a+bi}{c+di} \cdot \frac{c-di}{c-di} = \frac{(a+bi) \cdot (c-di)}{c^{2}+d^{2}}$
		\item $x = e^{Ln(x)} \leftrightarrow x = Ln(e^{x})$
	\end{itemize}
	\section{Ejercicio 1}
	Calcular parte real e imaginaria de los siguientes números complejos: \\
	a) $z_1 = \frac{3-2i}{1+4i} $
	b) $z_2 = \frac{1}{i} + \frac{1}{1+i}$
	c) $z_3 = cos(i)$
	d) $z_4 = sen(2+i)$
	\subsection{Solución}
	\[
	z_1 = \frac{3-2i}{1+4i} = \frac{3-2i}{1+4i} \cdot \frac{1-4i}{1-4i} = \frac{3-8-12i-2i}{17} = \boxed{-\frac{5}{17} - \frac{14}{17}i}
	\]
	\[
	z_2 = \frac{1}{i} + \frac{1}{1+i} = \frac{1+2i}{-1+i} = \frac{1+2i}{-1+i} \cdot \frac{-1-i}{-1-i} = \frac{-1+2-1i-2i}{2} = \boxed{\frac{1}{2} - \frac{3}{2}i}
	\]
	\[
	z_3 = cos(i) = \frac{e^{-1} + e}{2} = \frac{\frac{1}{e}+e}{2} = \boxed{\frac{1}{2e}+\frac{e}{2}}
	\]
	\[
	z_4 = sen(2+i) = \frac{e^{i(2+i)}-e^{-i(2+i)}}{2i} = \frac{e^{-1+2i}-e^{1-2i}}{2i} = -\frac{i}{2} (\frac{1}{e} \cdot (cos(2)+isen(2)) - e \cdot (cos(-2)+isen(-2)))
	\]
	\[
	= -\frac{i}{2} (\frac{1}{e} \cdot (cos(2)+isen(2)) - e \cdot (cos(2)-isen(2))) = \frac{1}{2e} (sen(2)-icos(2)) - \frac{e}{2} (-sen(2)-icos(2))
	\]
	\[
	= \boxed{(\frac{1}{2e}+\frac{e}{2})sen(2) + i(-\frac{1}{2e}+\frac{e}{2})cos(2)}
	\]
	\section{Ejercicio 2}
	Calcular el modulo y argumento de: \\
	a) $z_1 = 3^i$
	b) $z_2 = i^{i}$
	c) $z_3 = i^{3+i}$
	d) $z_4 = (1+i)^{2+i}$
	\subsection{Solución}
	\[
	z_1 = 3^{i} = e^{Ln(3^{i})} = e^{i Ln(3)} = \boxed{r=1 , \sigma = Ln (3)}
	\]
	\[
	z_2 = (e^{i\frac{\pi}{2}})^{i} = e^{-\frac{\pi}{2}} = \boxed{r=e^{-\frac{\pi}{2}}, \sigma = 0}
	\]
	\[
	z_3 = i^{3+i} = i^{3} \cdot i^{i} = (-i) \cdot e^{-\frac{\pi}{2}} = e^{-\frac{\pi}{2}} \cdot e^{i\frac{3\pi}{2}} = \boxed{r=e^{-\frac{\pi}{2}}, \sigma = \frac{\pi}{2}}
	\]
	\[
	z_4 = (1+i)^{2+i} = (1+i)^{2} \cdot (1+i)^{i} = (\sqrt{2} \cdot e^{i\frac{\pi}{4}})^{2} \cdot (\sqrt{2} \cdot e^{i\frac{\pi}{4}})^{i} = (2 \cdot e^{i\frac{\pi}{2}}) \cdot ((\sqrt{2})^{i} \cdot e^{-\frac{\pi}{4}})
	\]
	\[
	= 2e^{-\frac{\pi}{4}} \cdot e^{i\frac{\pi}{2}} \cdot e^{iLn(\sqrt{2})} = 2e^{-\frac{\pi}{4}} \cdot e^{i(\frac{\pi}{2} + Ln(\sqrt{2}))} = \boxed{r=2e^{-\frac{\pi}{4}}, \sigma = \frac{\pi}{2} + Ln(\sqrt{2})}
	\]
	\section{Ejercicio 3}
	Calcular las raíces cuartas de la unidad
	\subsection{Solución}
	\[
	z^{4} = 1 \rightarrow z = \sqrt[4]{1} e^{i(\frac{2k\pi}{4})} || k = 0,1,2,3
	\]
	\[
	k = 0 \rightarrow e^{0} = \boxed{1}
	\]
	\[
	k = 1 \rightarrow e^{i\frac{\pi}{2}} = \boxed{i}
	\]
	\[
	k = 2 \rightarrow e^{i\pi} = \boxed{-1}
	\]
	\[
	k = 3 \rightarrow e^{i3\frac{\pi}{2}} = \boxed{-i}
	\]
	\section{Ejercicio 4}
	Calcular las raíces cúbicas de -8
	\subsection{Solución}
	\[
	z^{3} = -8 \rightarrow \sqrt[3]{8} e^{i(\frac{\pi + 2k\pi}{3})} || k = 0,1,2
	\]
	\[
	k = 0 \rightarrow 2e^{i\frac{\pi}{3}} = 2(cos(\frac{\pi}{3}) + isen(\frac{\pi}{3})) = \boxed{1 + i\sqrt{3}} 
	\]
	\[
	k = 1 \rightarrow 2e^{i\pi} = \boxed{-2}
	\]
	\[
	k = 2 \rightarrow 2e^{i\frac{5\pi}{3}} = 2(cos(\frac{5\pi}{3}) + isen(\frac{5\pi}{3})) = \boxed{1 - i\sqrt{3}}
	\]
	\section{Ejercicio 5}
	Determina la forma binómica de $e^{\sqrt{i}}$
	\subsection{Solución}
	\[
	\sqrt{i} = \sqrt{e^{i\frac{\pi}{2}}} = e^{i(\frac{\frac{\pi}{2} + 2k\pi}{2})} || k = 0,1
	\]
	\[
	k = 0 \rightarrow e^{i\frac{\pi}{4}} = cos(\frac{\pi}{4}) + sen(\frac{\pi}{4}) = \frac{\sqrt{2}}{2} + i\frac{\sqrt{2}}{2}
	\]
	\[
	k = 1 \rightarrow e^{i\frac{5\pi}{4}} = cos(\frac{5\pi}{4}) + sen(5\frac{\pi}{4}) = -\frac{\sqrt{2}}{2} - i\frac{\sqrt{2}}{2}
	\]
	Opción 1
	\[
	e^{\sqrt{i}} = e^{\frac{\sqrt{2}}{2} + i\frac{\sqrt{2}}{2}} = e^{\frac{\sqrt{2}}{2}}(cos(\frac{\sqrt{2}}{2}) + i sen(\frac{\sqrt{2}}{2})) = \boxed{e^{\frac{\sqrt{2}}{2}}cos(\frac{\sqrt{2}}{2}) + ie^{\frac{\sqrt{2}}{2}}sen(\frac{\sqrt{2}}{2})}
	\]
	Opción 2
	\[
	e^{\sqrt{i}} = e^{\frac{-\sqrt{2}}{2} - i\frac{\sqrt{2}}{2}} = e^{\frac{-\sqrt{2}}{2}}(cos(-\frac{\sqrt{2}}{2}) + i sen(-\frac{\sqrt{2}}{2})) = \boxed{e^{-\frac{\sqrt{2}}{2}}cos(\frac{\sqrt{2}}{2}) - ie^{-\frac{\sqrt{2}}{2}}sen(\frac{\sqrt{2}}{2})}
	\]
	\section{Ejercicio 6}
	Expresa en forma binómica el número $(-\frac{\sqrt{3}}{2} + \frac{i}{2})^{6}$
	\subsection{Solución}
	\[
	(-\frac{\sqrt{3}}{2} + \frac{i}{2})^{6} = (e^{i\frac{5\pi}{6}})^{6} = e^{i5\pi} = e^{i\pi} = \boxed{-1}
	\]
	\section{Ejercicio 7}
	Calcula $\sqrt{-16-30i}$
	\subsection{Solución}
	\[
	\sqrt{-16-30i} = (x+iy) \rightarrow (\sqrt{-16-30i})^{2} = (x+iy)^{2} \rightarrow -16-30i = x^{2} + 2ixy - y^{2} \rightarrow
	\]
	\[
	x^{2} - y^{2} = -16
	\]
	\[
	2xy = -30 \rightarrow x = -\frac{15}{y}
	\]
	\[
	(-\frac{15}{y})^{2} - y^{2} = -16 \rightarrow \frac{225}{y^{2}} - y^{2} = -16 \rightarrow y^{4} +16y^{2} - 225 = 0
	\]
	\[
	\underrightarrow{y^{2} = t}
	\]
	\[
	t^{2} +16t -225 = 0 \rightarrow (t-25)(t+9) = 0 \rightarrow y=\pm5, x=\mp3
	\]
	\[
	\boxed{z_1 = -3 + 5i, z_2 = 3 -5i}
	\]
	\section{Ejercicio 8}
	Determinar $m \in \mathbb{R}$ de modo que $(2e^{i\sqrt{2}})^{m}$ sea un número real negativo
	\subsection{Solución}
	\[
	(2e^{i\sqrt{2}})^{m} = 2^{m} \cdot \underline{e^{im\sqrt{2}}}
	\]
	Nota: Nos fijamos en la parte imaginaria ya que $2^{m}$ siempre sera un valor positivo
	\[
	2^{m}(cos(m\sqrt{2}) + isen(m\sqrt{2}))
	\]
	Nota: Para que el número sea real negativo, el coseno debe ser negativo y el seno nulo. Por lo tanto, $\forall k \in \mathbb{Z} \rightarrow cos((2k+1)\pi) = -1, sen((2k+1)\pi) = 0$
	\[
	m\sqrt{2} = (2k+1)\pi \rightarrow \boxed{m = \frac{(2k+1)\pi}{\sqrt{2}} || k \in \mathbb{Z}}
	\]
	\section{Ejercicio 9}
	Determinar los números complejos no nulos tal que su quinta potencia, $z^{5}$, sea igual a su conjugado, es decir, $\overline{z}$
	\subsection{Solución}
	\[
	z = re^{i\sigma}, \overline{z} = z = re^{-i\sigma}
	\]
	\[
	(re^{i\sigma})^{5} = re^{-i\sigma} \rightarrow r^{5}e^{i\sigma5} = re^{-i\sigma}
	\]
	\[
	r^{5} = r \rightarrow r(r^{4} - 1) = 0 \rightarrow r = 0 X, r = -1 X, r = 1 \checkmark
	\]
	\[
	e^{i\sigma5} = e^{-i\sigma} \rightarrow 5\sigma = -\sigma + 2k\pi \rightarrow \sigma = \frac{k\pi}{3}
	\]
	\[
	k = 0 \rightarrow z = 1, \boxed{1^{5} = \overline{1}}
	\]
	\[
	k = 1 \rightarrow z = e^{i\frac{\pi}{3}}, \boxed{e^{i\frac{5\pi}{3}} = e^{-i\frac{\pi}{3}}}
	\]
	\[
	\boxed{\forall k \in \mathbb{N} \rightarrow (e^{i\frac{k\pi}{3}})^{5} = e^{-i\frac{k\pi}{3}}}
	\]
	\section{Ejercicio 10}
	Dado el polinomio $P(z) = z^{4} -6z^{3} + 24z^{2} -18z + 63$. Calcula $P(i\sqrt{3})$ y $P(-i\sqrt{3})$. Resuelve a continuación la ecuación $P(z) = 0$
	\subsection{Solución}
	\[
	P(i\sqrt{3}) = (i\sqrt{3})^{4} - 6(i\sqrt{3})^{3} + 24 (i\sqrt{3})^{2} - 18(i\sqrt{3}) + 63 = 9 + 18\sqrt{3}i - 72 - 18\sqrt{3}i + 63 = \boxed{0}
	\]
	\[
	\boxed{P(-i\sqrt{3}) = 0} 
	\]
	Nota: Si un número complejo es solución o raíz de un polinomio de coeficientes reales, su conjugado también lo es. \\
	\[
	P(z) = (z - i\sqrt{3})(z - (-i\sqrt{3})) \cdot Q(x) = (z^{2} + 3) \cdot Q(x) 
 	\]
 	\[
 	Q(x) = \frac{z^{4} -6z^{3} + 24z^{2} -18z + 63}{z^{2} + 3} = \text{...} = z^{2} - 6z + 21 = \text{...} = (z - (3+2\sqrt{3}i))(z - (3-2\sqrt{3}i))
 	\]
 	\[
 	P(z) = 0 \rightarrow \boxed{P(z) = (z - i\sqrt{3})(z - (-i\sqrt{3}))(z - (3+2\sqrt{3}i))(z - (3-2\sqrt{3}i))}
 	\]
 	\section{Ejercicio 11}
 	Dado el número complejo $z = -\sqrt{2+\sqrt{2}} + i\sqrt{2-\sqrt{2}}$, calcula $z^{2}$ en forma binómica. A continuación, expresa $z^{2}$ en forma exponencial y deduce la forma exponencial de $z$
 	\subsection{Solución}
 	 \[
 	 z^{2} = (-\sqrt{2+\sqrt{2}} + i\sqrt{2-\sqrt{2}})^{2} = 2 + \sqrt{2} - 2 + \sqrt{2} - 2i(\sqrt{2+\sqrt{2}})(\sqrt{2-\sqrt{2}}) = \boxed{2\sqrt{2} - 2i\sqrt{2}}
 	 \]
 	 \[
 	 \begin{cases}
 	 	r = \sqrt{(2\sqrt{2})^{2} + (2\sqrt{2})^{2}} = 4 \\
 	 	\sigma = arctg(\frac{-2\sqrt{2}}{2\sqrt{2}}) = -\frac{\pi}{4} = \frac{7\pi}{4}
 	 \end{cases}
 	 \]
 	 \[
 	 z^{2} = 4e^{\frac{7\pi}{4}}
 	 \]
 	 \[
 	 z = \sqrt{4e^{\frac{7\pi}{4}}} = 2e^{i\frac{\frac{7\pi}{4}+ 2k\pi}{2}} \text{ tal que } k = 0,1
 	 \]
 	 \[
 	 \begin{cases}
 	 	k = 0 & 2e^{i\frac{7\pi}{8}} \text{ Correcto, segundo cuadrante} \\
 	 	k = 1 & 2e^{i\frac{15\pi}{8}} \text{ Incorrecto, cuarto cuadrante}
 	 \end{cases}
 	 \]
 	 Nota: El número original estaba en el segundo cuadrante. Al elevar al cuadrado y después deshacerlo, podemos introducir soluciones irreales.


	\section{Ejercicio 12}
	Obtén la forma exponencial del número complejo $(1+i\sqrt{3})^{1-i}$.

	\subsection{Solución}

	Primero expresamos la base en forma exponencial

	\[
	1+i\sqrt{3}
	=
	2e^{i\frac{\pi}{3}}
	\]

	Separamos ahora el exponente:

	\[
	(1+i\sqrt{3})^{1-i}
	=
	2e^{i\frac{\pi}{3}}\left(2e^{i\frac{\pi}{3}}\right)^{-i}
	=
	2e^{i\frac{\pi}{3}} 2^{-i}\left(e^{i\frac{\pi}{3}}\right)^{-i}
	=
	2e^{i\frac{\pi}{3}} 2^{-i}e^{\frac{\pi}{3}}
	\]


	Para calcular el término $2^{-i}$ utilizamos la definición de potencia compleja:

	\[
	z^w=e^{w\ln(z)}
	\]

	\[
	2e^{i\frac{\pi}{3}} 2^{-i}e^{\frac{\pi}{3}}
	=
	2e^{\frac{\pi}{3}} e^{i\frac{\pi}{3}} e^{-i\ln(2)}
	=
	2e^{\frac{\pi}{3}} e^{i\left(\frac{\pi}{3}-\ln(2)\right)}
	\]

	\[
	\boxed{
	(1+i\sqrt{3})^{1-i}
	=
	2e^{\frac{\pi}{3}}
	e^{i\left(\frac{\pi}{3}-\ln(2)\right)}
	}
	\]
	
	\section{Ejercicio 13}
		Calcula la parte real e imaginaria del número complejo
		$z=\cos\left(\frac{\pi}{6}+2i\right)$.

		\subsection{Solución}

		Utilizamos la expresión:

		\[
		\cos(z)=\frac{e^{iz}+e^{-iz}}{2}
		\]

		\[
		\cos\left(\frac{\pi}{6}+2i\right)
		=
		\frac{
		e^{i\left(\frac{\pi}{6}+2i\right)}
		+
		e^{-i\left(\frac{\pi}{6}+2i\right)}
		}{2}
		=
		\frac{
		e^{-2}e^{i\frac{\pi}{6}}
		+
		e^2e^{-i\frac{\pi}{6}}
		}{2}
		\]

		Aplicando la fórmula de Euler:

		\[
		=
		\frac{1}{2}
		\left[
		e^{-2}
		\left(
		\cos\left(\frac{\pi}{6}\right)
		+i\sin\left(\frac{\pi}{6}\right)
		\right)
		+
		e^2
		\left(
		\cos\left(\frac{\pi}{6}\right)
		-i\sin\left(\frac{\pi}{6}\right)
		\right)
		\right]
		\]

		\[
		=
		\frac{1}{2}
		\left[
		e^{-2}\left(\frac{\sqrt{3}}{2}+\frac{i}{2}\right)
		+
		e^2\left(\frac{\sqrt{3}}{2}-\frac{i}{2}\right)
		\right]
		\]

		\[
		z
		=
		\frac{\sqrt{3}}{4}\left(e^2+e^{-2}\right)
		+
		i\frac{e^{-2}-e^2}{4}
		\]

		Por tanto:

		\[
		\boxed{
		\operatorname{Re}(z)
		=
		\frac{\sqrt{3}}{4}\left(e^2+e^{-2}\right)
		}
		\]

		\[
		\boxed{
		\operatorname{Im}(z)
		=
		\frac{e^{-2}-e^2}{4}
		}
		\]
	
	\section{Ejercicio 14}
		Obtén la forma binómica del número complejo

		\[
		z=
		\left(
		\frac{\sqrt{3}-i}{\sqrt{3}+i}
		\right)^4
		\left(
		\frac{1+i}{1-i}
		\right)^5.
		\]

		\subsection{Solución}

		Expresamos cada número complejo en forma exponencial:

		\[
		\sqrt{3}-i=2e^{-i\frac{\pi}{6}},
		\qquad
		\sqrt{3}+i=2e^{i\frac{\pi}{6}}
		\]

		\[
		1+i=\sqrt{2}e^{i\frac{\pi}{4}},
		\qquad
		1-i=\sqrt{2}e^{-i\frac{\pi}{4}}
		\]

		Por tanto:

		\[
		z=
		\left(
		\frac{2e^{-i\frac{\pi}{6}}}{2e^{i\frac{\pi}{6}}}
		\right)^4
		\left(
		\frac{\sqrt{2}e^{i\frac{\pi}{4}}}
		{\sqrt{2}e^{-i\frac{\pi}{4}}}
		\right)^5
		\]

		\[
		z=
		\left(e^{-i\frac{\pi}{3}}\right)^4
		\left(e^{i\frac{\pi}{2}}\right)^5
		=
		e^{-i\frac{4\pi}{3}}e^{i\frac{5\pi}{2}}
		=
		e^{i\frac{7\pi}{6}}
		\]

		Aplicando la fórmula de Euler:

		\[
		z=
		\cos\left(\frac{7\pi}{6}\right)
		+i\sin\left(\frac{7\pi}{6}\right)
		=
		-\frac{\sqrt{3}}{2}-\frac{i}{2}
		\]

		\[
		\boxed{
		z=-\frac{\sqrt{3}}{2}-\frac{1}{2}i
		}
		\]

	\section{Ejercicio 15}
		Resuelve la ecuación

		\[
		z^3+(-1+i)z^2+(1-i)z+i=0
		\]

		en el cuerpo de los números complejos, proporcionando las soluciones en forma binómica o exponencial.

		\subsection{Solución}

		Comprobamos que $z=-i$ es una raíz:

		\[
		(-i)^3+(-1+i)(-i)^2+(1-i)(-i)+i=0
		\]

		Por tanto, el polinomio es divisible entre $z+i$:

		\[
		z^3+(-1+i)z^2+(1-i)z+i
		=
		(z+i)(z^2-z+1)
		\]

		La ecuación queda:

		\[
		(z+i)(z^2-z+1)=0
		\]

		\[
		z=-i
		\]

		Para el segundo factor:

		\[
		z^2-z+1=0
		\]

		\[
		z
		=
		\frac{1\pm\sqrt{1-4}}{2}
		=
		\frac{1\pm i\sqrt{3}}{2}
		\]

		Por tanto, las soluciones en forma binómica son:

		\[
		\boxed{
		z_1=-i,
		\qquad
		z_2=\frac{1}{2}+\frac{\sqrt{3}}{2}i,
		\qquad
		z_3=\frac{1}{2}-\frac{\sqrt{3}}{2}i
		}
		\]

		En forma exponencial:

		\[
		\boxed{
		z_1=e^{-i\frac{\pi}{2}},
		\qquad
		z_2=e^{i\frac{\pi}{3}},
		\qquad
		z_3=e^{-i\frac{\pi}{3}}
		}
		\]
	
	\section{Ejercicio 16}
		Dado el número complejo

		\[
		z=\frac{1+i\sqrt{3}}{1-i\sqrt{3}},
		\]

		completa los siguientes apartados:

		\begin{enumerate}[label=\alph*)]
			\item Exprésalo en forma binómica y exponencial.
			\item Prueba que $z^4=z$.
			\item Determina las otras tres raíces de cuarto grado de $z$.
		\end{enumerate}

		\subsection{Solución}

		\subsubsection*{a)}

		\[
		z
		=
		\frac{1+i\sqrt{3}}{1-i\sqrt{3}}
		\frac{1+i\sqrt{3}}{1+i\sqrt{3}}
		=
		\frac{(1+i\sqrt{3})^2}{1+3}
		\]

		\[
		z
		=
		\frac{1+2i\sqrt{3}-3}{4}
		=
		-\frac{1}{2}+\frac{\sqrt{3}}{2}i
		\]

		\[
		\boxed{
		z=-\frac{1}{2}+\frac{\sqrt{3}}{2}i
		}
		\]

		Como $|z|=1$ y su argumento es $\frac{2\pi}{3}$:

		\[
		\boxed{
		z=e^{i\frac{2\pi}{3}}
		}
		\]

		\subsubsection*{b)}

		\[
		z^4
		=
		\left(e^{i\frac{2\pi}{3}}\right)^4
		=
		e^{i\frac{8\pi}{3}}
		=
		e^{i\left(2\pi+\frac{2\pi}{3}\right)}
		=
		e^{i\frac{2\pi}{3}}
		=
		z
		\]

		\[
		\boxed{
		z^4=z
		}
		\]

		\subsubsection*{c)}

		Las raíces cuartas de $z$ vienen dadas por:

		\[
		w_k
		=
		e^{i\frac{\frac{2\pi}{3}+2k\pi}{4}}
		=
		e^{i\left(\frac{\pi}{6}+\frac{k\pi}{2}\right)},
		\qquad
		k=0,1,2,3
		\]

		Para $k=0$:

		\[
		w_0
		=
		e^{i\frac{\pi}{6}}
		=
		\frac{\sqrt{3}}{2}+\frac{1}{2}i
		\]

		Para $k=1$:

		\[
		w_1
		=
		e^{i\frac{2\pi}{3}}
		=
		-\frac{1}{2}+\frac{\sqrt{3}}{2}i
		=
		z
		\]

		Para $k=2$:

		\[
		w_2
		=
		e^{i\frac{7\pi}{6}}
		=
		-\frac{\sqrt{3}}{2}-\frac{1}{2}i
		\]

		Para $k=3$:

		\[
		w_3
		=
		e^{i\frac{5\pi}{3}}
		=
		\frac{1}{2}-\frac{\sqrt{3}}{2}i
		\]

		Por tanto, las otras tres raíces cuartas de $z$ son:

		\[
		\boxed{
		e^{i\frac{\pi}{6}},
		\qquad
		e^{i\frac{7\pi}{6}},
		\qquad
		e^{i\frac{5\pi}{3}}
		}
		\]
	
	\section{Ejercicio 17}
		Resuelve la ecuación

		\[
		z^3+(5+3i)z^2+(4+7i)z=0
		\]

		en el cuerpo de los números complejos. Una vez obtenidas las soluciones, proporciona la forma binómica del inverso multiplicativo de cada solución.

		\subsection{Solución}

		Sacamos factor común $z$:

		\[
		z\left(z^2+(5+3i)z+(4+7i)\right)=0
		\]

		Por tanto:

		\[
		z_1=0
		\]

		Resolvemos la ecuación de segundo grado:

		\[
		z^2+(5+3i)z+(4+7i)=0
		\]

		\[
		z
		=
		\frac{-(5+3i)\pm\sqrt{(5+3i)^2-4(4+7i)}}{2}
		\]

		\[
		(5+3i)^2-4(4+7i)
		=
		16+30i-16-28i
		=
		2i
		\]

		Como

		\[
		(1+i)^2=2i,
		\]

		tenemos:

		\[
		z
		=
		\frac{-(5+3i)\pm(1+i)}{2}
		\]

		\[
		z_2
		=
		\frac{-5-3i+1+i}{2}
		=
		-2-i
		\]

		\[
		z_3
		=
		\frac{-5-3i-1-i}{2}
		=
		-3-2i
		\]

		Por tanto, las soluciones son:

		\[
		\boxed{
		z_1=0,
		\qquad
		z_2=-2-i,
		\qquad
		z_3=-3-2i
		}
		\]

		El número $z_1=0$ no tiene inverso multiplicativo.

		\[
		z_2^{-1}
		=
		\frac{1}{-2-i}
		=
		\frac{-2+i}{(-2)^2+1^2}
		=
		-\frac{2}{5}+\frac{1}{5}i
		\]

		\[
		z_3^{-1}
		=
		\frac{1}{-3-2i}
		=
		\frac{-3+2i}{(-3)^2+2^2}
		=
		-\frac{3}{13}+\frac{2}{13}i
		\]

		\[
		\boxed{
		z_2^{-1}
		=
		-\frac{2}{5}+\frac{1}{5}i,
		\qquad
		z_3^{-1}
		=
		-\frac{3}{13}+\frac{2}{13}i
		}
		\]
	
	\section{Ejercicio 18}
		Resuelve el siguiente sistema de ecuaciones en el cuerpo de los números complejos:

		\[
		\begin{cases}
		(3-i)x+(4+2i)y=2+6i,\\
		(4+2i)x-(2+3i)y=5+4i.
		\end{cases}
		\]

		\subsection{Solución}

		Aplicamos la regla de Cramer:

		\[
		\Delta
		=
		\begin{vmatrix}
		3-i & 4+2i\\
		4+2i & -(2+3i)
		\end{vmatrix}
		\]

		\[
		\Delta
		=
		(3-i)(-2-3i)-(4+2i)^2
		=
		(-9-7i)-(12+16i)
		=
		-21-23i
		\]

		\[
		\Delta_x
		=
		\begin{vmatrix}
		2+6i & 4+2i\\
		5+4i & -(2+3i)
		\end{vmatrix}
		\]

		\[
		\Delta_x
		=
		(2+6i)(-2-3i)-(4+2i)(5+4i)
		=
		(14-18i)-(12+26i)
		=
		2-44i
		\]

		\[
		x
		=
		\frac{\Delta_x}{\Delta}
		=
		\frac{2-44i}{-21-23i}
		=
		1+i
		\]

		\[
		\Delta_y
		=
		\begin{vmatrix}
		3-i & 2+6i\\
		4+2i & 5+4i
		\end{vmatrix}
		\]

		\[
		\Delta_y
		=
		(3-i)(5+4i)-(2+6i)(4+2i)
		=
		(19+7i)-(-4+28i)
		=
		23-21i
		\]

		\[
		y
		=
		\frac{\Delta_y}{\Delta}
		=
		\frac{23-21i}{-21-23i}
		=
		i
		\]

		\[
		\boxed{
		x=1+i,
		\qquad
		y=i
		}
		\]
	
	\section{Ejercicio 19}
		Resuelve la ecuación compleja

		\[
		\sin(z)=5.
		\]

		\subsection{Solución}

		Utilizamos la expresión:

		\[
		\sin(z)=\frac{e^{iz}-e^{-iz}}{2i}
		\]

		\[
		\frac{e^{iz}-e^{-iz}}{2i}=5
		\]

		Multiplicando por $2ie^{iz}$:

		\[
		e^{2iz}-1=10ie^{iz}
		\]

		\[
		e^{2iz}-10ie^{iz}-1=0
		\]

		Hacemos el cambio:

		\[
		w=e^{iz}
		\]

		\[
		w^2-10iw-1=0
		\]

		\[
		w
		=
		\frac{10i\pm\sqrt{(10i)^2+4}}{2}
		=
		\frac{10i\pm\sqrt{-96}}{2}
		=
		i(5\pm2\sqrt{6})
		\]

		Por tanto:

		\[
		e^{iz}=i(5+2\sqrt{6})
		\]

		o bien:

		\[
		e^{iz}=i(5-2\sqrt{6})
		\]

		Como

		\[
		i=e^{i\left(\frac{\pi}{2}+2k\pi\right)},
		\qquad k\in\mathbb{Z},
		\]

		obtenemos:

		\[
		iz
		=
		\ln(5+2\sqrt{6})
		+
		i\left(\frac{\pi}{2}+2k\pi\right)
		\]

		\[
		z
		=
		\frac{\pi}{2}+2k\pi
		-i\ln(5+2\sqrt{6})
		\]

		Para la segunda solución, utilizamos:

		\[
		5-2\sqrt{6}
		=
		\frac{1}{5+2\sqrt{6}}
		\]

		\[
		\ln(5-2\sqrt{6})
		=
		-\ln(5+2\sqrt{6})
		\]

		\[
		z
		=
		\frac{\pi}{2}+2k\pi
		+i\ln(5+2\sqrt{6})
		\]

		Por tanto:

		\[
		\boxed{
		z
		=
		\frac{\pi}{2}+2k\pi
		\pm i\ln(5+2\sqrt{6}),
		\qquad k\in\mathbb{Z}
		}
		\]

	\section{Ejercicio 20}
		Resuelve la ecuación compleja

		\[
		\cos(z)+i\sqrt{3}=0.
		\]

		\subsection{Solución}

		\[
		\cos(z)=-i\sqrt{3}
		\]

		Utilizamos la expresión:

		\[
		\cos(z)=\frac{e^{iz}+e^{-iz}}{2}
		\]

		\[
		\frac{e^{iz}+e^{-iz}}{2}=-i\sqrt{3}
		\]

		Multiplicando por $2e^{iz}$:

		\[
		e^{2iz}+1=-2i\sqrt{3}e^{iz}
		\]

		\[
		e^{2iz}+2i\sqrt{3}e^{iz}+1=0
		\]

		Hacemos el cambio:

		\[
		w=e^{iz}
		\]

		\[
		w^2+2i\sqrt{3}w+1=0
		\]

		\[
		w
		=
		\frac{-2i\sqrt{3}\pm\sqrt{(2i\sqrt{3})^2-4}}{2}
		=
		\frac{-2i\sqrt{3}\pm4i}{2}
		\]

		\[
		w_1=i(2-\sqrt{3}),
		\qquad
		w_2=-i(2+\sqrt{3})
		\]

		Para la primera solución:

		\[
		e^{iz}
		=
		(2-\sqrt{3})e^{i\left(\frac{\pi}{2}+2k\pi\right)}
		\]

		\[
		iz
		=
		\ln(2-\sqrt{3})
		+
		i\left(\frac{\pi}{2}+2k\pi\right)
		\]

		\[
		z
		=
		\frac{\pi}{2}+2k\pi
		-i\ln(2-\sqrt{3})
		\]

		Como

		\[
		2-\sqrt{3}=\frac{1}{2+\sqrt{3}},
		\]

		queda:

		\[
		z
		=
		\frac{\pi}{2}+2k\pi
		+i\ln(2+\sqrt{3})
		\]

		Para la segunda solución:

		\[
		e^{iz}
		=
		(2+\sqrt{3})e^{i\left(-\frac{\pi}{2}+2k\pi\right)}
		\]

		\[
		iz
		=
		\ln(2+\sqrt{3})
		+
		i\left(-\frac{\pi}{2}+2k\pi\right)
		\]

		\[
		z
		=
		-\frac{\pi}{2}+2k\pi
		-i\ln(2+\sqrt{3})
		\]

		Por tanto:

		\[
		\boxed{
		z
		=
		2k\pi
		\pm
		\left(
		\frac{\pi}{2}
		+i\ln(2+\sqrt{3})
		\right),
		\qquad k\in\mathbb{Z}
		}
		\]
	
	\section{Ejercicio 21}
		Determina el valor de los parámetros $a$ y $b$ para que el número complejo

		\[
		z=\frac{3b-2ai}{4-3i}
		\]

		sea real y de módulo unidad.

		\subsection{Solución}

		Multiplicamos por el conjugado del denominador:

		\[
		z
		=
		\frac{3b-2ai}{4-3i}
		\frac{4+3i}{4+3i}
		=
		\frac{(3b-2ai)(4+3i)}{25}
		\]

		\[
		z
		=
		\frac{12b+6a+(9b-8a)i}{25}
		\]

		Para que $z$ sea real, su parte imaginaria debe ser nula:

		\[
		9b-8a=0
		\qquad\Longrightarrow\qquad
		a=\frac{9}{8}b
		\]

		Como $z$ es real y tiene módulo unidad:

		\[
		z=\pm1
		\]

		\[
		\frac{12b+6a}{25}=\pm1
		\]

		Sustituyendo $a=\frac{9}{8}b$:

		\[
		\frac{12b+6\left(\frac{9}{8}b\right)}{25}
		=
		\frac{3b}{4}
		=
		\pm1
		\]

		\[
		b=\pm\frac{4}{3}
		\]

		\[
		a=\frac{9}{8}b
		=
		\pm\frac{3}{2}
		\]

		Por tanto:

		\[
		\boxed{
		(a,b)=
		\left(\frac{3}{2},\frac{4}{3}\right)
		\quad\text{o}\quad
		(a,b)=
		\left(-\frac{3}{2},-\frac{4}{3}\right)
		}
		\]
	
	\section{Ejercicio 22}
		Determina de forma analítica el objeto geométrico que forman los números complejos que satisfacen la relación

		\[
		|z-2i|=2|z+3|.
		\]

		\subsection{Solución}

		Sea

		\[
		z=x+iy.
		\]

		Entonces:

		\[
		|z-2i|
		=
		|x+i(y-2)|
		=
		\sqrt{x^2+(y-2)^2}
		\]

		\[
		|z+3|
		=
		|(x+3)+iy|
		=
		\sqrt{(x+3)^2+y^2}
		\]

		Por tanto:

		\[
		\sqrt{x^2+(y-2)^2}
		=
		2\sqrt{(x+3)^2+y^2}
		\]

		Elevando al cuadrado:

		\[
		x^2+(y-2)^2
		=
		4\left((x+3)^2+y^2\right)
		\]

		\[
		x^2+y^2-4y+4
		=
		4x^2+24x+36+4y^2
		\]

		\[
		3x^2+3y^2+24x+4y+32=0
		\]

		Dividiendo entre $3$ y completando cuadrados:

		\[
		x^2+8x+y^2+\frac{4}{3}y+\frac{32}{3}=0
		\]

		\[
		(x+4)^2+\left(y+\frac{2}{3}\right)^2
		=
		\frac{52}{9}
		\]

		Por tanto, el lugar geométrico es una circunferencia de centro

		\[
		\left(-4,-\frac{2}{3}\right)
		\]

		y radio

		\[
		\frac{\sqrt{52}}{3}
		=
		\frac{2\sqrt{13}}{3}.
		\]

		\[
		\boxed{
		\left(x+4\right)^2+
		\left(y+\frac{2}{3}\right)^2
		=
		\left(\frac{2\sqrt{13}}{3}\right)^2
		}
		\]
	
	\section{Ejercicio 23}
		Demuestra que el lugar geométrico de los números complejos $z$ que satisfacen la ecuación

		\[
		|z-b|+|z+b|=2a
		\]

		es una elipse, donde $a,b\in\mathbb{R}$.

		\subsection{Solución}

		Sea

		\[
		z=x+iy.
		\]

		Entonces:

		\[
		|z-b|
		=
		\sqrt{(x-b)^2+y^2},
		\qquad
		|z+b|
		=
		\sqrt{(x+b)^2+y^2}
		\]

		Por tanto:

		\[
		\sqrt{(x-b)^2+y^2}
		+
		\sqrt{(x+b)^2+y^2}
		=
		2a
		\]

		Aislamos uno de los radicales:

		\[
		\sqrt{(x-b)^2+y^2}
		=
		2a-\sqrt{(x+b)^2+y^2}
		\]

		Elevando al cuadrado:

		\[
		(x-b)^2+y^2
		=
		4a^2+(x+b)^2+y^2
		-
		4a\sqrt{(x+b)^2+y^2}
		\]

		\[
		a\sqrt{(x+b)^2+y^2}
		=
		a^2+bx
		\]

		Volvemos a elevar al cuadrado:

		\[
		a^2\left((x+b)^2+y^2\right)
		=
	(a^2+bx)^2
		\]

		\[
		a^2x^2+2a^2bx+a^2b^2+a^2y^2
		=
		a^4+2a^2bx+b^2x^2
		\]

		\[
		(a^2-b^2)x^2+a^2y^2
		=
		a^2(a^2-b^2)
		\]

		Dividiendo entre $a^2(a^2-b^2)$:

		\[
		\boxed{
		\frac{x^2}{a^2}
		+
		\frac{y^2}{a^2-b^2}
		=
		1
		}
		\]

		Esta es la ecuación de una elipse centrada en el origen, con semiejes

		\[
		a
		\qquad\text{y}\qquad
		\sqrt{a^2-b^2},
		\]

		y focos situados en

		\[
		(-b,0)
		\qquad\text{y}\qquad
		(b,0).
		\]

		Para que la elipse sea no degenerada debe cumplirse:

		\[
		\boxed{a>|b|}
		\]
	
	\section{Ejercicio 24}
		Determina los números complejos que cumplen cada una de las siguientes condiciones, proporcionando además una interpretación geométrica de los conjuntos de soluciones:

		\begin{enumerate}[label=\alph*)]
			\item $\operatorname{Re}(iz)-z\overline{z}=0$.
			\item $\operatorname{Im}\left((3+i)z\right)=2$.
		\end{enumerate}

		\subsection{Solución}

		Sea

		\[
		z=x+iy,
		\qquad x,y\in\mathbb{R}.
		\]

		\subsubsection*{a)}

		\[
		iz=i(x+iy)=-y+ix
		\]

		\[
		\operatorname{Re}(iz)=-y,
		\qquad
		z\overline{z}=|z|^2=x^2+y^2
		\]

		Por tanto:

		\[
		-y-x^2-y^2=0
		\]

		\[
		x^2+y^2+y=0
		\]

		Completando cuadrados:

		\[
		x^2+\left(y+\frac{1}{2}\right)^2
		=
		\frac{1}{4}
		\]

		El conjunto de soluciones es:

		\[
		\boxed{
		\left\{
		z=x+iy\in\mathbb{C}
		\;\middle|\;
		x^2+\left(y+\frac{1}{2}\right)^2=\frac{1}{4}
		\right\}
		}
		\]

		Geométricamente, es una circunferencia de centro

		\[
		\left(0,-\frac{1}{2}\right)
		\]

		y radio

		\[
		\frac{1}{2}.
		\]

		\subsubsection*{b)}

		\[
		(3+i)z
		=
		(3+i)(x+iy)
		=
		3x-y+i(x+3y)
		\]

		Por tanto:

		\[
		\operatorname{Im}\left((3+i)z\right)
		=
		x+3y
		\]

		\[
		x+3y=2
		\]

		El conjunto de soluciones es:

		\[
		\boxed{
		\left\{
		z=x+iy\in\mathbb{C}
		\;\middle|\;
		x+3y=2
		\right\}
		}
		\]

		Geométricamente, es la recta:

		\[
		\boxed{
		y=-\frac{1}{3}x+\frac{2}{3}
		}
		\]

	\section{Ejercicio 25}
		Determina el conjunto de números complejos que cumplen por separado las siguientes condiciones. En el caso de que los afijos tengan una forma geométrica reconocible, indica de cuál se trata.

		\begin{enumerate}[label=\alph*)]
			\item
			\[
			\arg\left(\frac{z-1}{z+1}\right)=\frac{\pi}{4}.
			\]

			\item
			\[
			|z|=z+2\overline{z}+1+2i.
			\]
		\end{enumerate}

		\subsection{Solución}

		Sea

		\[
		z=x+iy,
		\qquad x,y\in\mathbb{R}.
		\]

		\subsubsection*{a)}

		\[
		\frac{z-1}{z+1}
		=
		\frac{x-1+iy}{x+1+iy}
		\frac{x+1-iy}{x+1-iy}
		\]

		\[
		\frac{z-1}{z+1}
		=
		\frac{x^2+y^2-1+2iy}{(x+1)^2+y^2}
		\]

		Para que su argumento sea $\frac{\pi}{4}$, las partes real e imaginaria deben ser iguales y positivas:

		\[
		x^2+y^2-1=2y,
		\qquad y>0
		\]

		\[
		x^2+(y-1)^2=2,
		\qquad y>0
		\]

		Por tanto:

		\[
		\boxed{
		\left\{
		z=x+iy\in\mathbb{C}
		\;\middle|\;
		x^2+(y-1)^2=2,\ y>0
		\right\}
		}
		\]

		Geométricamente, es la parte situada en el semiplano superior de la circunferencia de centro $(0,1)$ y radio $\sqrt{2}$.

		\subsubsection*{b)}

		\[
		|z|
		=
		x+iy+2(x-iy)+1+2i
		\]

		\[
		\sqrt{x^2+y^2}
		=
		3x+1+i(2-y)
		\]

		Como el miembro izquierdo es real:

		\[
		2-y=0
		\qquad\Longrightarrow\qquad
		y=2
		\]

		\[
		\sqrt{x^2+4}=3x+1
		\]

		Elevando al cuadrado:

		\[
		x^2+4=(3x+1)^2
		\]

		\[
		8x^2+6x-3=0
		\]

		\[
		x
		=
		\frac{-6\pm\sqrt{36+96}}{16}
		=
		\frac{-3\pm\sqrt{33}}{8}
		\]

		Debe cumplirse $3x+1\geq 0$, por lo que únicamente es válida la solución:

		\[
		x=\frac{-3+\sqrt{33}}{8}
		\]

		Por tanto:

		\[
		\boxed{
		z=\frac{\sqrt{33}-3}{8}+2i
		}
		\]

		Geométricamente, el conjunto de soluciones está formado por un único punto del plano complejo.

	\section{Ejercicio 26}
		Sabiendo que

		\[
		\tanh(z)=\frac{e^z-e^{-z}}{e^z+e^{-z}},
		\]

		proporciona de forma razonada dos números complejos $z_1$ y $z_2$, con partes real e imaginaria no nulas, que cumplan

		\[
		|\tanh(z)|=1.
		\]

		\subsection{Solución}

		\[
		\left|
		\frac{e^z-e^{-z}}{e^z+e^{-z}}
		\right|
		=1
		\]

		\[
		\left|
		\frac{e^{2z}-1}{e^{2z}+1}
		\right|
		=1
		\]

		\[
		|e^{2z}-1|
		=
		|e^{2z}+1|
		\]

		Sea $z=x+iy$. Entonces:

		\[
		e^{2z}
		=
		e^{2x}e^{2iy}
		=
		e^{2x}\bigl(\cos(2y)+i\sin(2y)\bigr)
		\]

		La igualdad

		\[
		|e^{2z}-1|=|e^{2z}+1|
		\]

		indica que $e^{2z}$ debe estar a la misma distancia de $1$ y de $-1$, por lo que su parte real debe ser nula:

		\[
		\operatorname{Re}(e^{2z})
		=
		e^{2x}\cos(2y)
		=
		0
		\]

		Como $e^{2x}\neq 0$:

		\[
		\cos(2y)=0
		\]

		\[
		2y=\frac{\pi}{2}+k\pi
		\qquad\Longrightarrow\qquad
		y=\frac{\pi}{4}+\frac{k\pi}{2}
		\]

		Tomando valores no nulos para $x$ e $y$, podemos elegir:

		\[
		\boxed{
		z_1=1+\frac{\pi}{4}i,
		\qquad
		z_2=2+\frac{3\pi}{4}i
		}
		\]

		Ambos cumplen $|\tanh(z)|=1$ y tienen partes real e imaginaria no nulas.
\end{document}